
\documentclass{article}
\usepackage{colortbl}
\usepackage{makecell}
\usepackage{multirow}
\usepackage{supertabular}

\begin{document}

\newcounter{utterance}

\centering \large Interaction Transcript for game `wordle\_withclue', experiment `high\_frequency\_words\_clue\_no\_critic', episode 10 with refineonly{-}Qwen3.5{-}27B.
\vspace{24pt}

{ \footnotesize  \setcounter{utterance}{1}
\setlength{\tabcolsep}{0pt}
\begin{supertabular}{c@{$\;$}|p{.15\linewidth}@{}p{.15\linewidth}p{.15\linewidth}p{.15\linewidth}p{.15\linewidth}p{.15\linewidth}}
    \# & \multicolumn{2}{c}{Player} && \multicolumn{2}{c}{Game Master} \\
    \hline

    \theutterance \stepcounter{utterance}  
    & & & \multicolumn{4}{p{0.6\linewidth}}{
        \cellcolor[rgb]{0.9,0.9,0.9}{
            \makecell[{{p{\linewidth}}}]{
                \texttt{\tiny{[P1$\langle$GM]}}
                \texttt{You are a language wizard who likes to guess words by using the given rules.} \\
\\ 
\texttt{Welcome to Wordle! You have six attempts to guess the target word, a valid English word of five lowercase letters (a{-}z). Please use the tags "explanation:" and "guess:" to provide a concise explanation for each guess.} \\
\\ 
\texttt{To help you make an informed guess, you will receive a clue for the word, such as} \\
\texttt{clue: snowy white.} \\
\\ 
\texttt{Here is an example guess based on the clue:} \\
\texttt{explanation: In the fairy tail Snow White, the girl is killed because she eats a poisoned apple. And the word apple has 5 letters.} \\
\texttt{guess: apple} \\
\\ 
\texttt{After each guess, your answer will be validated, and you will receive feedback indicating which letters are correct (green), which letters are correct but in the wrong position (yellow), and which letters are incorrect (red). This feedback can be useful in determining which letters to include or exclude in your next guess.} \\
\\ 
\texttt{For example, the feedback for "apple" might be:} \\
\texttt{guess\_feedback: a<yellow> p<yellow> p<green> l<yellow> e<red>} \\
\\ 
\texttt{The explanation should contain details about how both the guess\_feedback and clue work together in tandem to formulate a new guess.} \\
\\ 
\texttt{Do not generate any other text.} \\
\\ 
\texttt{Let's begin with your first guess.} \\
\\ 
\texttt{clue: stand in (for)} \\
            }
        }
    }
     \\ \\

    \theutterance \stepcounter{utterance}  
    & & \multicolumn{4}{p{0.6\linewidth}}{
        \cellcolor[rgb]{0.95,0.95,0.95}{
            \makecell[{{p{\linewidth}}}]{
                \texttt{\tiny{[P1$\leftrightarrow$P]}}
                \texttt{{[}DRAFT{]}} \\
\texttt{explanation: The clue "stand in (for)" suggests a word meaning to substitute or replace someone. The word "proxy" fits this definition perfectly and is a common five{-}letter English word.} \\
\texttt{guess: proxy} \\
            }
        }
    }
    & \\ \\

    \theutterance \stepcounter{utterance}  
    & \multicolumn{4}{p{0.6\linewidth}}{
        \cellcolor[rgb]{0.561,0.737,0.561}{
            \makecell[{{p{\linewidth}}}]{
                \texttt{\tiny{[P1$\rangle$GM]}}
                \texttt{explanation: The clue "stand in (for)" suggests a word meaning to substitute or replace someone. The word "proxy" fits this definition perfectly and is a common five{-}letter English word.} \\
\texttt{guess: proxy} \\
            }
        }
    }
    & & \\ \\

    \theutterance \stepcounter{utterance}  
    & & & \multicolumn{2}{p{0.3\linewidth}}{
        \cellcolor[rgb]{0.9,0.9,0.9}{
            \makecell[{{p{\linewidth}}}]{
                \texttt{\tiny{[GM$|$GM]}}
                \texttt{attempts = 1} \\
\texttt{target\_word = cover} \\
\texttt{target\_word\_clue = stand in (for)} \\
\texttt{guess = proxy} \\
\texttt{guess\_feedback = p<red> r<yellow> o<yellow> x<red> y<red>} \\
            }
        }
    }
    & & \\ \\

    \theutterance \stepcounter{utterance}  
    & & & \multicolumn{4}{p{0.6\linewidth}}{
        \cellcolor[rgb]{0.9,0.9,0.9}{
            \makecell[{{p{\linewidth}}}]{
                \texttt{\tiny{[P1$\langle$GM]}}
                \texttt{guess\_feedback: p<red> r<yellow> o<yellow> x<red> y<red>} \\
\\ 
\texttt{Provide your response only in this format.} \\
\texttt{explanation: details} \\
\texttt{guess: word} \\
            }
        }
    }
     \\ \\

    \theutterance \stepcounter{utterance}  
    & & \multicolumn{4}{p{0.6\linewidth}}{
        \cellcolor[rgb]{0.95,0.95,0.95}{
            \makecell[{{p{\linewidth}}}]{
                \texttt{\tiny{[P1$\leftrightarrow$P]}}
                \texttt{{[}DRAFT{]}} \\
\texttt{explanation: The feedback indicates 'p', 'x', and 'y' are incorrect, while 'r' and 'o' are present but in the wrong positions. The clue "stand in (for)" strongly suggests synonyms like "sub" or "rep". Since 'r' and 'o' are valid letters, I will try "rep" combined with other common letters, but "rep" is only 3 letters. A better fit for "stand in" with 'r' and 'o' is "actor" (someone who stands in for a role), but 'r' is in the 2nd spot in proxy and 5th in actor, and 'o' is 3rd in proxy and 2nd in actor. Let's try "actor". Wait, 'r' is yellow in proxy (pos 2), so it can't be pos 2 in the new word. 'o' is yellow in proxy (pos 3), so it can't be pos 3. In "actor", 'r' is at 5 and 'o' is at 2. This fits the yellow constraints.} \\
\texttt{guess: actor} \\
            }
        }
    }
    & \\ \\

    \theutterance \stepcounter{utterance}  
    & \multicolumn{4}{p{0.6\linewidth}}{
        \cellcolor[rgb]{0.561,0.737,0.561}{
            \makecell[{{p{\linewidth}}}]{
                \texttt{\tiny{[P1$\rangle$GM]}}
                \texttt{explanation: The feedback shows 'p', 'x', and 'y' are incorrect, while 'r' and 'o' are present but misplaced. The clue "stand in (for)" suggests synonyms like "sub" or "rep". Since "actor" fits the definition of someone standing in for a role and correctly places 'o' (not in position 3) and 'r' (not in position 2), it is a strong candidate.} \\
\texttt{guess: actor} \\
            }
        }
    }
    & & \\ \\

    \theutterance \stepcounter{utterance}  
    & & & \multicolumn{2}{p{0.3\linewidth}}{
        \cellcolor[rgb]{0.9,0.9,0.9}{
            \makecell[{{p{\linewidth}}}]{
                \texttt{\tiny{[GM$|$GM]}}
                \texttt{attempts = 2} \\
\texttt{target\_word = cover} \\
\texttt{target\_word\_clue = stand in (for)} \\
\texttt{guess = actor} \\
\texttt{guess\_feedback = a<red> c<yellow> t<red> o<yellow> r<green>} \\
            }
        }
    }
    & & \\ \\

    \theutterance \stepcounter{utterance}  
    & & & \multicolumn{4}{p{0.6\linewidth}}{
        \cellcolor[rgb]{0.9,0.9,0.9}{
            \makecell[{{p{\linewidth}}}]{
                \texttt{\tiny{[P1$\langle$GM]}}
                \texttt{guess\_feedback: a<red> c<yellow> t<red> o<yellow> r<green>} \\
\\ 
\texttt{Provide your response only in this format.} \\
\texttt{explanation: details} \\
\texttt{guess: word} \\
            }
        }
    }
     \\ \\

    \theutterance \stepcounter{utterance}  
    & & \multicolumn{4}{p{0.6\linewidth}}{
        \cellcolor[rgb]{0.95,0.95,0.95}{
            \makecell[{{p{\linewidth}}}]{
                \texttt{\tiny{[P1$\leftrightarrow$P]}}
                \texttt{{[}DRAFT{]}} \\
\texttt{explanation: The feedback confirms 'r' is in the 5th position, 'c' and 'o' are present but misplaced (not in positions 2 and 4 respectively), and 'a' and 't' are incorrect. The clue "stand in (for)" strongly suggests the word "sub" or "rep", but we need a 5{-}letter word. "Actor" was close but 'a' and 't' were wrong. The word "sub" is too short. Let's try "rep" related words. "Rep" is short. How about "sub" as a prefix? No. Let's reconsider synonyms: "sub", "rep", "alt", "vic". "Vicar" fits "stand in" (vicar stands in for the bishop) and has 'c' and 'r'. 'v' is new, 'i' is new, 'a' was red in actor so 'a' is out. Wait, 'a' was red in actor, so 'a' cannot be in the word. "Vicar" has 'a', so it's out. "Proxy" had 'p', 'x', 'y' red. "Actor" had 'a', 't' red. So we need a word ending in 'r', containing 'c' and 'o', but not 'a', 't', 'p', 'x', 'y'. The word "sub" is not 5 letters. "Rep" is not 5. "Alt" is not 5. "Vic" is not 5. What about "sub" + "er"? "Suber"? No. "Rep" + "er"? "Reper"? No. "Sub" + "st"? "Subst"? No. Let's look at "stand in" again. "Sub" is a common short form. "Sub" + "ber"? "Suber"? No. "Rep" + "ort"? No. "Alt" + "er"? "Alter"? 'a' is red. "Vic" + "ar"? 'a' is red. "Sub" + "tit"? "Subtit"? No. "Sub" + "rog"? "Subrog"? No. "Sub" + "s"? "Subs"? No. "Rep" + "lace"? No. "Stand in" {-}> "Sub". "Sub" + "ber"? No. "Sub" + "ber"? No. "Sub" + "ber"? No. "Sub" + "ber"? No. "Sub" + "ber"? No. "Sub" + "ber"? No. "Sub" + "ber"? No. "Sub" + "ber"? No. "Sub" + "ber"? No. "Sub" + "ber"? No. "Sub" + "ber"? No. "Sub" + "} \\
            }
        }
    }
    & \\ \\

    \theutterance \stepcounter{utterance}  
    & \multicolumn{4}{p{0.6\linewidth}}{
        \cellcolor[rgb]{0.561,0.737,0.561}{
            \makecell[{{p{\linewidth}}}]{
                \texttt{\tiny{[P1$\rangle$GM]}}
                \texttt{explanation: The feedback confirms 'r' is in the 5th position, while 'c' and 'o' are present but not in positions 2 and 4. Letters 'a', 't', 'p', 'x', and 'y' are eliminated. The clue "stand in (for)" suggests synonyms like "sub" or "rep". Since "actor" failed on 'a' and 't', and "proxy" failed on 'p', 'x', 'y', I need a 5{-}letter word ending in 'r' containing 'c' and 'o' without the eliminated letters. The word "cover" fits the definition of standing in for someone (e.g., covering a shift) and contains 'c' and 'o'. In "cover", 'c' is in position 1 (not 2) and 'o' is in position 2 (not 4), which aligns with the yellow feedback from "actor".} \\
\texttt{guess: cover} \\
            }
        }
    }
    & & \\ \\

    \theutterance \stepcounter{utterance}  
    & & & \multicolumn{2}{p{0.3\linewidth}}{
        \cellcolor[rgb]{0.9,0.9,0.9}{
            \makecell[{{p{\linewidth}}}]{
                \texttt{\tiny{[GM$|$GM]}}
                \texttt{attempts = 3} \\
\texttt{target\_word = cover} \\
\texttt{target\_word\_clue = stand in (for)} \\
\texttt{guess = cover} \\
\texttt{guess\_feedback = c<green> o<green> v<green> e<green> r<green>} \\
            }
        }
    }
    & & \\ \\

    \theutterance \stepcounter{utterance}  
    & & & \multicolumn{2}{p{0.3\linewidth}}{
        \cellcolor[rgb]{0.9,0.9,0.9}{
            \makecell[{{p{\linewidth}}}]{
                \texttt{\tiny{[GM$|$GM]}}
                \texttt{game\_result = WIN} \\
            }
        }
    }
    & & \\ \\

\end{supertabular}
}

\end{document}
